\begin{ex}
Một chiếc xe chở 3 tạ gạo và 5 yến ngô. Chiếc xe chở tất cả bao nhiêu ki-lô-gam?
\shortans{350}
\loigiai{
$300 + 50 = 350$ kg.
}
\end{ex}

\begin{ex}
Một tấm biển quảng cáo hình chữ nhật dài 4 m, rộng 2 m. Diện tích tấm biển là bao nhiêu mét vuông?
\shortans{8}
\loigiai{
$4 \times 2 = 8\text{ m}^2$.
}
\end{ex}

\begin{ex}
Bạn Lan rửa bát hết 4 phút 15 giây. Tính theo giây bạn Lan rửa hết bao nhiêu giây?
\shortans{255}
\loigiai{
$4 \times 60 + 15 = 255$ giây.
}
\end{ex}

\begin{ex}
Một bể bơi mini có diện tích đáy là $25\text{ m}^2$. Đổi sang đề-xi-mét vuông là bao nhiêu?
\shortans{2500}
\loigiai{
$25 \times 100 = 2500\text{ dm}^2$.
}
\end{ex}

\begin{ex}
Mỗi bao xi măng nặng 50 kg. Một tấn xi măng có bao nhiêu bao?
\shortans{20}
\loigiai{
$1000 : 50 = 20$ bao.
}
\end{ex}

\begin{ex}
Năm 1428 Lê Lợi lên ngôi vua sau khởi nghĩa Lam Sơn. Năm đó thuộc thế kỉ thứ mấy? (Ghi số tự nhiên).
\shortans{15}
\loigiai{
Thế kỉ XV (15).
}
\end{ex}

\begin{ex}
Năm 1789 vua Quang Trung đại phá quân Thanh. Năm đó thuộc thế kỉ thứ mấy? (Ghi số tự nhiên).
\shortans{18}
\loigiai{
Thế kỉ XVIII (18).
}
\end{ex}

\begin{ex}
Một mặt bàn học hình chữ nhật dài 12 dm, rộng 6 dm. Diện tích mặt bàn là bao nhiêu đề-xi-mét vuông?
\shortans{72}
\loigiai{
$12 \times 6 = 72\text{ dm}^2$.
}
\end{ex}

\begin{ex}
Tính: 3 tấn + 25 tạ = \dots tạ.
\shortans{55}
\loigiai{
$30 + 25 = 55$ tạ.
}
\end{ex}

\begin{ex}
Tính: 5 tạ - 18 yến = \dots yến.
\shortans{32}
\loigiai{
$50 - 18 = 32$ yến.
}
\end{ex}

\begin{ex}
Tính: $2\text{ m}^2\ 50\text{ dm}^2 + 3\text{ m}^2\ 50\text{ dm}^2 = \dots\text{ m}^2$.
\shortans{6}
\loigiai{
$250 + 350 = 600\text{ dm}^2 = 6\text{ m}^2$.
}
\end{ex}

\begin{ex}
Đổi: 2 phút 40 giây = \dots giây.
\shortans{160}
\loigiai{
$120 + 40 = 160$ giây.
}
\end{ex}

\begin{ex}
Đổi: 420 giây = \dots phút.
\shortans{7}
\loigiai{
$420 : 60 = 7$ phút.
}
\end{ex}

\begin{ex}
Một đội trồng rừng tuần thứ nhất trồng được 3 tấn cây giống, tuần thứ hai trồng được 25 tạ. Cả hai tuần trồng được bao nhiêu tạ cây giống?
\shortans{55}
\loigiai{
$30 + 25 = 55$ tạ.
}
\end{ex}

\begin{ex}
Căn phòng hình chữ nhật dài 8 m, rộng 5 m. Người ta dùng gạch vuông cạnh 40 cm để lát nền. Diện tích căn phòng là bao nhiêu mét vuông?
\shortans{40}
\loigiai{
$8 \times 5 = 40\text{ m}^2$.
}
\end{ex}

\begin{ex}
Đổi: $40\text{ m}^2 = \dots\text{ dm}^2$.
\shortans{4000}
\loigiai{
$40 \times 100 = 4000\text{ dm}^2$.
}
\end{ex}

\begin{ex}
Chiến thắng Điện Biên Phủ diễn ra năm 1954. Năm đó thuộc thế kỉ thứ mấy? (Ghi số tự nhiên).
\shortans{20}
\loigiai{
Thế kỉ XX (20).
}
\end{ex}

\begin{ex}
Cách mạng Tháng Tám thành công năm 1945. Năm 1945 thuộc thế kỉ thứ mấy?
\shortans{20}
\loigiai{
Thế kỉ XX (20).
}
\end{ex}

\begin{ex}
Một con chim di cư bay liên tục trong 1 phút 20 giây. Thời gian bay là bao nhiêu giây?
\shortans{80}
\loigiai{
$60 + 20 = 80$ giây.
}
\end{ex}

\begin{ex}
Một con trâu nặng 400 kg. Hỏi 5 con trâu như thế nặng bao nhiêu tấn?
\shortans{2}
\loigiai{
$400 \times 5 = 2000\text{ kg} = 2\text{ tấn}$.
}
\end{ex}

\begin{ex}
Một thửa ruộng thu hoạch được 2 tấn 4 tạ thóc. Đổi ra ki-lô-gam được bao nhiêu?
\shortans{2400}
\loigiai{
$2000 + 400 = 2400$ kg.
}
\end{ex}

\begin{ex}
Một viên gạch có diện tích $400\text{ cm}^2$. Đổi ra đề-xi-mét vuông là bao nhiêu?
\shortans{4}
\loigiai{
$400 : 100 = 4\text{ dm}^2$.
}
\end{ex}

\begin{ex}
Tính: $\frac{1}{4}$ thế kỉ = \dots năm.
\shortans{25}
\loigiai{
$100 : 4 = 25$ năm.
}
\end{ex}

\begin{ex}
Tính: $\frac{1}{3}$ phút = \dots giây.
\shortans{20}
\loigiai{
$60 : 3 = 20$ giây.
}
\end{ex}

\begin{ex}
Đổi: 4 tạ 5 yến = \dots kg.
\shortans{450}
\loigiai{
$400 + 50 = 450$ kg.
}
\end{ex}
